\section{Finite-Sample and Nonbinary Extensions}\label{app:extensions}
\subsection{When the Address Gap Becomes Visible}\label{app:addressfinite}
Consider the address contrast with $D$ data coordinates, either the binary construction of Theorem~\ref{thm:universal} or the $s$-ary construction below. The data coordinates take values in $\{-1,1\}$. The fixed score $H=DYZ_J\mathbf1\{J=A\}$ has $\E H=\Delta$ and $\E H^2=D$ under every population. With $\gamma=\sqrt n/(1+\sqrt n)$, its shrunken sample mean has risk
\[
 \frac{\gamma^2D}{n}+\left((1-\gamma)^2-\frac{\gamma^2}{n}\right)\Delta^2
   =\frac{D}{(1+\sqrt n)^2}.
\]
The zero estimator also has risk at most one. Branching recovers the sign $Yd$ and has exact minimax risk $(1+\sqrt n)^{-2}$ by Appendix~\ref{app:sign}.

For a converse take independent uniform address coordinates and independent fair data bits, with profile law $\mu$. On the subfamily $P_\theta(x,y)=\mu(x)(1+\theta y d(x))/2$, the gain is $\theta$. For a fixed mask $T$, put $u_T=\E_\mu[d\mid X_T]$. A transcript $(x_T,y)$ has density proportional to $1+\theta y u_T(x_T)$, so its Fisher information is
\[
 I_T(\theta)=\E_\mu\!\left[\frac{u_T^2}{1-\theta^2u_T^2}\right]
 \le\frac{1}{D(1-\eta^2)},\qquad |\theta|\le\eta<1.
\]
Here $|u_T|\le1$ and the address projection calculation gives $\E u_T^2\le1/D$. The bound is uniform in the chosen mask. Conditional scores are martingale differences, so choosing masks from earlier respondents cannot increase total information beyond $n/[D(1-\eta^2)]$. The same bound holds if the current label selects the mask: for either selected mask the law of $u_T$ is symmetric under flipping all data bits, and $\E[u_T^2/(1\pm\theta u_T)]=\E[u_T^2/(1-\theta^2u_T^2)]$.

The prior $\eta^{-1}\cos^2(\pi\theta/(2\eta))$ on $[-\eta,\eta]$ has Fisher information $\pi^2/\eta^2$. The van Trees argument in Appendix~\ref{app:general} therefore gives
\[
 R_{\rm fixed}(b,n)\ge
 \frac{1}{n/[D(1-\eta^2)]+\pi^2/\eta^2}.
\]
Minimizing the denominator over $\eta^2\in(0,1)$ gives $(\sqrt{n/D}+\pi)^2$. Consequently
\[
 \frac{D}{(\sqrt n+\pi\sqrt D)^2}
 \le R_{\rm fixed}(b,n)
 \le\min\left\{1,\frac{D}{(1+\sqrt n)^2}\right\}.
\]
This proves Proposition~\ref{prop:finiteaddress}. The ratio is at least $D/2$ whenever $n\ge60D$, since $\sqrt{60}/(\sqrt{60}+\pi)>1/\sqrt2$. This threshold is sufficient, not necessary; the result does not identify the exact finite-sample fixed-form risk.

\subsection{Features with More Than Two Answers}\label{app:alphabet}
Let each of the $m$ coordinates have at most $s\ge2$ possible answers. Pad a depth-$b$ tree to $b$ distinct coordinates as in Appendix~\ref{app:universal}. Its first $b-1$ answers determine the entire query set, giving at most $s^{b-1}$ sets. Every set contains the root coordinate, so there are also at most $\binom{m-1}{b-1}$. The same simulation and the same conditional second-moment argument give Corollary~\ref{cor:alphabet} and its finite-program counterpart.

For sharpness put $\kappa=b-1$ and $D=s^\kappa$. Use $\kappa$ independent uniform $s$-valued address coordinates and $D$ independent fair data bits; features with fewer than $s$ possible answers are permitted. The contrast is $d=Z_A$. A mask with $u$ address coordinates and $r$ data coordinates has
\[
 \big\|\E[d\mid X_T]\big\|_2^2=r s^{u-2\kappa}\le s^{-\kappa},\qquad u+r\le \kappa+1,
\]
because $r s^{1-r}\le1$ for integers $r\ge1$. Thus the variance dual gives $V_{\rm fixed}\ge D$. Querying the complete address and one independent uniform data coordinate yields a pointwise unbiased score with second moment $D$, while branching evaluates $d$ with second moment one. The coordinate count is $\kappa+D$. If a common alphabet of exactly $s$ symbols is required, map each data symbol to a sign and use a lower-bound population supported on two symbols of opposite signs. The bounds allow arbitrary population laws.

Both complete menus span sums of functions of at most $b$ coordinates: each adaptive leaf indicator depends on at most $b$ coordinates, and fixed forms span every such function. Thus the identification spaces still coincide. The binary Fourier-degree language is replaced by interaction order in a product basis.

\subsection{Approaching an Identification Floor}\label{app:floor-rate}
For any $d$, define $\delta=\inf_{v\in\Hh_\A}\|d-v\|_\infty$. The indistinguishable pair in Appendix~\ref{app:id} proves $R_\A(d,n)\ge\delta^2$ for every $n$. For each $v\in\Hh_\A$, use a design with conditional score mean $Yv(X)$ and maximal second moment approaching $V(v,\A)$. Relative to the original target $Yd(X)$, its mean has absolute bias at most $\|d-v\|_\infty$. Hence
\begin{equation}
 \delta^2\le R_\A(d,n)
 \le\inf_{v\in\Hh_\A}\left\{\|d-v\|_\infty^2+\frac{V(v,\A)}n\right\}.
 \label{eq:floor-envelope}
\end{equation}
A best uniform approximant exists in this finite-dimensional subspace and has finite $V$. For any such $v_\star$, the excess risk is at most $V(v_\star,\A)/n$. On a common binary cap, the two menus have the same $\delta$; applying Theorem~\ref{thm:universal} to each approximant compares their variance terms by $K_{m,b}$. This is an upper bound on approach to the floor, not a matching excess-risk rate.

For example, on twenty binary features let $d(x)=0.8x_1+0.2\prod_{j=1}^{20}x_j$ and set $b=5$. Its degree exceeds the cap, but $v=0.8x_1$ has uniform approximation error $0.2$. The uniform signed witness proportional to $\prod_{j=1}^{20}x_j$ annihilates every degree-five function, proving that this error is optimal. Both menus satisfy
\[
 0.04\le R_\A(d,n)\le0.04+0.64/n.
\]
For a general dense logistic update, determining an accurate approximant can itself be difficult. More respondents alone cannot remove a positive identification floor. Additional structural assumptions or a larger per-person cap are required for consistency.
